\vspace{-4pt}
\section{\lhb{} as a Pipeline for Loop Engineering Evaluation in Long Horizon Coding}\label{sec:design}

\rev{Figure~\ref{fig:overview} summarizes the pipeline components, while the numbered stages below specify their construction order.} \lhb{} constructs loop engineering evaluation instances in which dependency structure controls test release, completed work remains live as regression obligations, and each run records loop trace metrics for planning, routing, verification, state continuity, and obligation retention. \textbf{(1) Task Collection} gathers source artifacts from three authentic real world sources. \textbf{(2) Task Preprocessing} normalizes the source artifacts into atomic candidate tasks. \textbf{(3) Task Selection} filters candidates by two source agnostic thresholds.
\rev{\textbf{(4) Development Unit Definition and Intra Task Relation Recovery} defines the units and recovers the dependency DAG that governs ready frontier release.}
\rev{\textbf{(5) Task Instrumentation} then attaches the public instruction, executable environment, and per unit test obligations using the recovered graph.}

\vspace{-3pt}
\subsection{Task Collection}\label{sec:task-collection}

% Overview / pipeline figure: three task sources -> uniform sandbox.
\begin{figure}[!t]
  \centering
  \includegraphics[width=\textwidth]{figures/assets/overview.pdf}
  \caption{Overview of the \lhb{} construction and evaluation pipeline.}
  \label{fig:overview}
\end{figure}


We collect source artifacts from three authentic real world sources spanning a broad spectrum of coding domains, task shapes, and dependency depths. This source mix varies routing burden and obligation retention pressure while keeping all tasks under the same loop engineering evaluation contract. \textbf{(1) Course Labs.} We collect 112 university programming projects from the official course homepages of 30 universities and from public GitHub mirrors of their assignment releases, retrieving each as its complete native bundle of handout, skeleton code, and instructor side reference tests. \textbf{(2) PR Sequences.} We pull the full commit history and merged pull request stream of 56 actively maintained GitHub repositories, all under continuous community development with high star counts and long maintenance windows. \textbf{(3) Research Evolutions.} We gather 212 papers from 48 venues together with the arXiv tex of their cited prior work and forward citers, restricted to papers that are simultaneously highly cited and accompanied by a fully open source official implementation under an OSI approved license. A reference filter retains the 95 papers connected by load bearing inheritance edges.
The three sources jointly anchor the benchmark in authentic developer, educator, and researcher artifacts and yield complementary long horizon structures for evaluating coding agent loops. Per source coverage appears in Appendix~\ref{app:per-task-stats}.

\vspace{-3pt}
\subsection{Task Preprocessing}\label{sec:preprocessing}

We normalize the collected artifacts into atomic candidate tasks before selection so later stages can represent heterogeneous sources through a uniform evaluation contract.

\parabf{Course Labs.} \ding{182}~unify handouts written in mixed languages into a single English version; \ding{183}~reorganize scattered skeleton files under one task root. Each project becomes one candidate task, yielding 112 candidates across the collected projects.

\parabf{PR Sequences.} \ding{182}~clone each repository with full git history and snapshot the first commit as the base codebase; \ding{183}~collect its merged pull requests and extract per pull request diffs split by role (e.g., test, source, auxiliary); \ding{184}~run a Claude Code based extraction step~\cite{claude_code} that performs repository level exploration and emits aggregated segments validated against a fixed acceptance standard. The procedure yields 148 candidate tasks across the 56 repositories.

\parabf{Research Evolutions.} \ding{182}~select a high impact early seed paper in the collected venue set and collect the arXiv tex of the seed together with its cited prior work and forward citers; \ding{183}~run a Claude Code based classification step that labels each reference as a shallow related work mention or a load bearing inheritance edge. \ding{184}~absorb every neighbor connected by a load bearing edge into the cluster and repeat the classification on the newly added papers until the closure admits no further additions. Each closed cluster becomes one candidate task, yielding 47 candidates across the 95 retained papers; 26 survive Selection.
Across all three sources, heterogeneous native test runners are wrapped behind a uniform task entry point, ensuring that every evaluated loop observes the same runtime contract while the source specific dependency evidence is preserved.
Appendix~\ref{app:example-task} presents one materialized \lhb{} task, from candidate recovery to the loop trace metrics recorded during evaluation.

\vspace{-3pt}
\subsection{Task Selection}\label{sec:dedup-quality}

Candidate tasks pass two source agnostic thresholds that retain workloads of sufficient scale for routing burden, state continuity, and regression obligation retention to become observable. \parabf{Temporal Span} is at least 2.5 months, measured as the handout declared project duration for Course Labs, the merge time interval between the first and last commit of the region for PR Sequences, and the publication date span between the earliest and latest paper in the evolution for Research Evolutions. \parabf{Solution Scale} is at least 1{,}200 in the source specific scale. Each source has its own natural patch length and base codebase size. Per source values are reported in Appendix~\ref{app:per-task-stats}.
After filtering, 112 tasks remain as the released loop engineering evaluation suite.

\vspace{-3pt}
\subsection{Intra Task Relation Recovery}\label{sec:dag-framework}

\rev{A materialized development unit is a source grounded, separately testable acceptance unit $u=(r_u,s_u,p_u,\Delta_u,T_u)$, where $r_u$ is its requirement, $s_u$ its file or symbol scope, $p_u$ its prerequisite set, $\Delta_u$ its reference patch contribution, and $T_u$ its standard tests. Unit boundaries follow source provenance: merged PRs or evidence supported PR segments for PR Sequences, released modules or milestones for Course Labs, and load bearing methodological increments for Research Evolutions. Changes that cannot be separated without breaking the acceptance contract are assigned to the nearest prerequisite unit or merged rather than split artificially.}

We formalize each task as a directed acyclic graph over its development units, allowing loop engineering evaluation to observe how a loop routes work through unit level prerequisites. \rev{The fixed DAG is an evaluation contract rather than a claim that real development is always monotonic or interface stable. An evaluated loop may revise earlier implementations provided that the completed predecessor obligations remain satisfied; tasks that intentionally replace earlier requirements through exploratory redesign fall outside this abstraction.} The four admitted patterns are \emph{sequential PR chains} along the merged commit history, \emph{structural module reuse} when $v$ edits a file or symbol that $u$ introduced, \emph{functional producer consumer API edges} when $v$ calls or imports an authoritative definition created by $u$, and \emph{compositional layering} when $v$ extends a subclass, schema, or interface declared by $u$.

% Temporal Dependency Types / DAG Construction Figure
\begin{center}
  \begin{minipage}{\textwidth}
    \centering
    \includegraphics[width=\linewidth]{figures/assets/dependency-types.pdf}
    \captionof{figure}{Four prerequisite patterns in a materialized task.}
    \label{fig:dependency-types}
  \end{minipage}
\end{center}


An edge is admitted only under unambiguous evidence on the gold artifacts. The first kind of evidence is hunk overlap on the same file. The second kind of evidence is symbol level dependence through an import, call, or subclass of an authoritative definition introduced by $u$. Spurious chains from hot monoliths, generated artifacts, and lockfiles are excluded by a denylist applied before edge admission. A file touched by most PRs in a repository therefore contributes no edges, and a regenerated lockfile never anchors a successor. The full admission rules appear in Appendix~\ref{app:task-schema}. On $G$, the \emph{ready frontier} at checkpoint $t$ is
\[
R_t \;=\; \{\, v \in V \setminus C_t \;:\; \forall (u,v) \in E,\; u \in C_t \,\},
\]
where $C_t$ collects nodes whose tests all pass at $t$.
For loop engineering evaluation, edges encode how local requirements become persistent obligations that remain active as later units are attempted.
The tests for each unit are evaluated in a workspace where predecessor tests remain enforced, measuring progress on later units together with preservation of earlier obligations.

\vspace{-3pt}
\subsection{Task Instrumentation}\label{sec:description-refine}

We turn each selected candidate into a self contained evaluation instance whose evaluation contract contains a formalized task instruction, a reproducible environment, and discriminative tests. The gold solution stays outside the evaluated loop context and is used only to validate executable obligations and compute reference statistics before evaluation begins.

\ding{182}~\textbf{Instruction Formalization.} We treat the gold diff against the base codebase as ground truth and recover the developer's acceptance intent from it, since PR titles, course handouts, and paper method sections rarely specify the full executable contract on their own.
A Claude Code based recovery step~\cite{claude_code} backed by Claude Opus 4.7~\cite{anthropic_opus47} produces a single acceptance instruction while keeping the implementation path hidden.
The recovery step has two roles in the evaluation contract. It removes evidence of the implementation path, including concrete edit locations and gold diff structure. It retains the public acceptance contract that the test suite binds to, including public signatures and expected output formats.
A symbol enters the instruction when a test assertion depends on naming it.
The same recovery step strips repository names, course identifiers, paper titles, and author attributions to reduce prompt level source recognition.

\ding{183}~\textbf{Environment and Test Suite Synthesis.}
A Claude Code based environment synthesis step traverses the source frame from upstream preprocessing, collects runtime dependencies, and writes a \texttt{Dockerfile} for the environment, with a companion \texttt{docker-compose.yaml} when the task spans multiple services.
\rev{After development units and their dependency DAG have been recovered,} the materialization step traverses the DAG in topological order so each unit test is authored against the prerequisite state used for ready frontier release.
At each unit $u$, the materialization pipeline drafts missing tests on top of the current environment, applies $u$'s gold patch on the recorded base, and runs the new tests under a three trial fail to pass check that admits $T(u)$ when the full gold solution flips every test to pass (\emph{solvability}), the empty workspace flips none (\emph{non triviality}), and the gold solution restricted to $u$'s DAG ancestors still leaves at least one test failing (\emph{discriminativeness} with respect to $u$'s own contribution).
A test with no execution result is attributed to the environment and triggers a \texttt{Dockerfile} revision, while a test that executes yet remains nonpassing under the gold patch is attributed to the assertion and rewritten in place.
After every unit has been visited, a closure pass replays the full gold patch against the accumulated suite on the recorded base. This grounds loop traces in executable obligations and excludes textual milestones from the scoring contract.
Units with unrepaired environment or suite issues after the computational budget are discarded.

\vspace{-3pt}
\subsection{Evaluation Protocol}\label{sec:protocol}

As illustrated in Figure~\ref{fig:protocol}, the evaluated loop receives the task instruction and per unit requirement bundles as plain files, while gold solutions and the ready unit identity remain withheld. \rev{Checkpoint outcomes and the active obligation state are retained for evaluator side scoring; the evaluated loop can inspect the attached tests, use project native tests, and author tests of its own. Appendix~\ref{app:task-schema} gives the complete visibility contract in Table~\ref{tab:visibility-contract}.} \rev{The evaluation runtime records the visitation sequence in the loop trace and compares the prerequisite structure in the recorded plan with the source recovered DAG. It accepts any valid topological execution order rather than requiring reproduction of the historical total order.}

\parabf{Topologically Gated Test Release.}
The evaluation runtime releases tests layer by layer along the dependency DAG, with every multi predecessor node acting as a gate.
\rev{Here, release and sealing refer only to evaluator side activation for scoring, not disclosure of the attached test files or bindings.}
A gate node $v$ joins the ready frontier after all predecessors lie in $C_t$, and the tests of its descendants stay sealed until $T(v)$ flips from fail to pass.
The gate affects scoring rather than editing permission. The evaluated loop remains free to edit across layers, and the recorded loop trace reports whether its work satisfies the dependency frontier without being told which unit is currently ready.
Once a unit clears the gate, its tests and the tests of its predecessors are kept enforced as regression tests on every later layer, scoring subsequent edits against both new work and completed obligations.
\rev{Regression Rate therefore measures preservation of previously satisfied behavior without staged evaluator feedback about checkpoint outcomes or the active obligation state.}

\parabf{Dual Container Snapshot Pipeline.}
We separate edit execution from test adjudication across two containers, allowing the evaluated loop to continue editing while the evaluation runtime independently samples workspace state and updates released obligations.
Container A hosts the working tree and is the sole write target for the evaluated loop.
Container B contains the full test suite and the reference environment, and reads from A through snapshots managed by the evaluation runtime.
On a fixed time cadence, the evaluation runtime samples A and diffs the new snapshot against the queue head. Empty diffs are dropped, while non empty diffs are appended to the queue and shipped into B for the currently released layer to run against.
Snapshots form a queue of monotonic edit states, allowing \lhb{} to compute loop trace metrics at real change points as well as at the final state.
